\section{Introduction}
\label{sec:introduction}


Multimodal deep search has emerged as a critical direction for multimodal large language models (MLLMs), enabling them to evolve from passive visual understanding systems into agents that actively search evidence, verify facts, and reason over knowledge-intensive visual queries \citep{huang2026vision,feng2026gen,chen2026unify}.
However, frontier multimodal search agents remain difficult to reproduce, as their training data, code are often proprietary or insufficiently disclosed \citep{seed2.0,huang2026vision,singh2025openai, kimiteam2026kimik25visualagentic}. As a result, the community still lacks a fully open recipe for building, analyzing, and improving strong multimodal search agents.

Among these missing components, high-quality training data is a central bottleneck. The strongest frontier systems are still largely dominated by well-funded commercial corporations \citep{Claude_4,comanici2025gemini}, where the data sources, filtering criteria, expert demonstrations, and tool-use trajectories are typically kept private. 
This makes it difficult to reproduce advanced multimodal search capabilities or systematically study which data properties are essential for agentic search behavior.
The issue is even more pronounced in multimodal settings, where effective training data must capture image-grounded understanding, multi-hop retrieval, evidence verification, and long-horizon tool use rather than simple visual question answering. 
Therefore, releasing high-quality training data is crucial for making frontier multimodal search agent research more transparent, reproducible, and accessible.

Beyond data, training multimodal search agents also poses unique challenges, especially when applying agentic reinforcement learning (agentic RL) \citep{fan2026exploring,webwatcher,huang2026vision} to long-horizon tool-use settings. 
Agentic search trajectories involve multiple rounds of reasoning, tool invocation, and observation integration, where a single malformed call, timeout, irrelevant query, or repeated failure can invalidate the remaining rollout. 
Simply discarding such trajectories wastes useful pre-failure reasoning, while training on the full rollout introduces noisy gradients from meaningless post-failure tokens. 
Another practical challenge is that real-world visual inputs are often imperfect, such as blurred photos, low-resolution thumbnails, skewed documents, and crowded screenshots. 
In these cases, searching alone is insufficient, and the agent must first crop, enhance, rectify, or parse the visual evidence before reliable search can begin. 
However, most existing multimodal search agents focus mainly on retrieval and do not jointly address robust visual pre-processing and failure-aware long-horizon RL.

In this work, we introduce \textbf{OpenSearch-VL}, a fully open recipe for training frontier multimodal deep search agents with agentic RL. 
Our recipe addresses the above challenges from data, tools, and training. 
First, we develop a dedicated data curation pipeline to build high-quality training data. 
Starting from the Wikipedia hyperlink graph, we sample multi-hop entity paths and convert them into multi-hop VQA instances by rewriting intermediate entities into fuzzy descriptions, followed by a carefully designed filtering mechanism.
This design avoids single-hop image lookup shortcuts and encourages the agent to learn multi-hop search and reasoning behaviors.
This pipeline yields two training datasets \textbf{SearchVL-SFT-36k} for SFT and \textbf{SearchVL-RL-8k} for agentic RL.
Second, we build a tool environment that goes beyond retrieval-only multimodal agent. 
In addition to search, the agent is equipped with OCR, cropping, sharpening, super-resolution, and perspective correction, allowing it to handle imperfect visual inputs in real-world scenarios before querying external knowledge. 
Finally, we develop an agentic RL algorithm based on GRPO ~\citep{guo2025deepseek} for long-horizon multimodal tool use, where multi-step interactions often lead to cascading tool failures. To address this issue, we introduce fatal-aware token masking that removes invalid post-failure suffixes from optimization, while preserving useful pre-failure reasoning through one-sided advantage clamping. This enables the model to learn from partially successful trajectories without being affected by noisy gradients from failed rollouts.
Together, these designs enable OpenSearch-VL to learn robust long-horizon search behavior over multimodal evidence in real-world scenarios.


Experiments across multimodal deep search benchmarks show that OpenSearch-VL consistently improves over strong baselines. 
For example, compared with the Qwen3-VL-30B-A3B \citep{Qwen3-VL} agentic baseline, our model improves the average score from 47.8 to 61.6, with large gains on VDR (+13.3) \citep{vdr}, MMSearch (+24.5) \citep{mmsearch}, FVQA (+10.2) \citep{wang2017fvqa}, and InfoSeek (+16.2) \citep{chen2023can}. 
Moreover, OpenSearch-VL achieves comparable or even better performance than proprietary commercial models on several benchmarks.
In summary, our main contributions can be summarized as follows:

\begin{itemize}[leftmargin=*, noitemsep]
    \item We introduce \textbf{OpenSearch-VL}, a fully open recipe for training frontier multimodal deep search agents. 
    We will release the training data, code, and models to provide an open foundation for reproducible research on multimodal agentic search.
    \item We build the key components required for training advanced multimodal search agents, including high-quality image-grounded multi-hop training data, a diverse tool environment, and a multi-turn fatal-aware GRPO algorithm. 
    \item Extensive experiments demonstrate the effectiveness of our recipe. For example, our trained OpenSearch-VL-30B-A3B brings an average improvement of 13.8 points across 7 multimodal deep search benchmarks.
\end{itemize}

\vspace{-0.1in}